\section{Preliminaries}
\label{sec:prelim}

\paragraph{Notation.}
For multivariate forecasting with $d$ variables, the task at origin $t$ is to
predict the next $H$ steps from the most recent $L$ observations. Let
$\xctx_t\in\R^{L\times d}$ be the observed context, $\ytrue_t\in\R^{H\times d}$
its realized future, and $f:\R^{L\times d}\to\R^{H\times d}$ a forecaster with
prediction $\ybase_t=f(\xctx_t)$. We assume $f$ has already been trained and
deployed: we may query it, but its parameters are never updated. The quantity this work centres on is the \emph{model-relative residual}
$\resid_t = \ytrue_t - \ybase_t$, defined jointly by the data and the particular
forecaster and observable only after the whole target window has elapsed. A
retrieval-augmented forecaster additionally has a memory $\mathcal{M}_t$ of past
items and emits a corrected forecast $\ycorr_t$; methods differ in whether those
items hold prior contexts with the futures that followed or, as we propose, the
same forecaster's prior residuals. Throughout, \emph{causal} means
free of lookahead: since $\resid_i$ needs the complete future at origin $i$, an
item may enter $\mathcal{M}_t$ only if $i+H<t$.

\begin{problem}[Retrieval-augmented correction of a trained forecast]
\label{prob:main}
Given a trained forecaster $f$, a query context $\xctx_t$ with base forecast
$\ybase_t=f(\xctx_t)$, and a memory $\mathcal{M}_t$ restricted to items whose
targets were fully observed before $t$, construct an additive correction $c_t$
so that $\ycorr_t=\ybase_t+c_t$ is more accurate than $\ybase_t$ with respect to
$\ytrue_t$, without updating $f$ and without observing any part of $\ytrue_t$.
\end{problem}

Problem~\ref{prob:main} isolates the question of
Section~\ref{sec:introduction}: the base forecast, the data, and the information
boundary are fixed, so any accuracy difference is attributable to what
$\mathcal{M}_t$ stores and how it is used.
